\begin{textsample}{POS Dim 5 – vem\_ed\_19 – Score 4.00 – vem\_ed\_19/t088.txt}  \label{ex:f5_pos_131}
Que conselho você daria para professores interessados em manter Intercâmbios Virtuais por um longo tempo?
Tentamos fazer os professores enxergarem oportunidades que os projetos COIL oferecem para além da duração das colaborações.
Se é uma parceria nova, encorajamos a oferecer o mesmo projeto por vários semestres.
O planejamento da primeira edição leva mais tempo, mas da segunda oferta em diante, são mudanças menores.
Utilizar a mesma estrutura de projeto com outros grupos de estudantes, e com parceiros em diversos países, traz múltiplas perspectivas.
Outro aspecto que vejo crescer é a pesquisa: professores trabalhando juntos investigam a metodologia de Intercâmbios Virtuais em seu campo de conhecimento e em suas turmas.
E publicam artigos em periódicos acadêmicos sobre essa inovação.
Se têm parceiros em vários países, podem realizar pesquisas multidisciplinares em diversas regiões do mundo.
Há muita demanda para pesquisa sobre Intercâmbios Virtuais: sobre o desenvolvimento de habilidades dos alunos, sobre a aprendizagem sob as perspectivas comparadas das turmas e como isso muda de acordo com o país — os resultados são parecidos ou diferentes?
Há tantas oportunidades para produzir artigos, periódicos ou até mesmo capítulos de livros.
Qual sua visão para os Intercâmbios Virtuais nos próximos cinco anos?
Talvez por causa de meu papel, vejo todos os benefícios que os Intercâmbios Virtuais podem trazer.
Acho que COIL veio para ficar e vai crescer à medida que mais instituições se juntarem a essa comunidade.
A pandemia ajudou de certa forma as instituições que buscam oferecer experiências de internacionalização para seus estudantes.
Parte de minha visão e esperança é que haja mais apoio estável para iniciativas de Intercâmbio Virtual nas instituições.
Muitas começam sem equipe, apenas com uma pessoa que \textbf{vê} COIL como uma grande oportunidade para seus professores, estudantes e para a instituição como um todo.
Depois de um tempo, \textbf{conseguem} mais recursos.
É importante reconhecer o papel das estruturas de suporte nas instituições - tecnológico, de design de projetos, de parcerias, linguístico.
Acho que o desenvolvimento da tecnologia de Inteligência Artificial para auxiliar nas traduções simultâneas deve facilitar o trabalho colaborativo, especialmente o de pequenos grupos de estudantes.
Também preciso mencionar as \textbf{redes} de coordenadores de Intercâmbios Virtuais nas diferentes instituições.
É importante aprender uns com os outros, porque há um conhecimento incrível nas diferentes regiões do mundo, e manter vivas essas \textbf{redes}.
É uma comunidade maravilhosa, que pode se beneficiar de recursos e ferramentas que ajudem a todas as instituições.

% matched lemmas: conseguir, rede, ver
\end{textsample}
